\newpage
\section{Conclusão}
\label{sec:conclusao}

Este trabalho caracterizou a frota eletrificada brasileira quanto às químicas de
bateria e às trajetórias de eletrificação até 2040 e quantificou, por meio de uma
Análise de Fluxo de Materiais e Energia, os fluxos associados à fabricação de
células LFP. Os principais resultados, organizados por objetivo específico, são
os seguintes.

\paragraph{Linha de base das químicas (Objetivo~1).}
Na amostra de mais de 509~mil veículos eletrificados emplacados, predomina o LFP
(48,12\%), seguido do NMC (32,75\%) e do NiMH (18,38\%), em participações por
número de veículos. O LFP domina entre os BEV — nove dos doze modelos de maior
volume —, e o NiMH está associado aos híbridos convencionais da Toyota. Como a
massa de bateria por veículo varia em uma ordem de grandeza entre categorias, essa
linha de base precisa ser convertida em massa ou em capacidade antes de alimentar
uma análise de fluxo de materiais da frota: com o \textit{mix} de vendas do
PDE~2035, os BEV respondem por pouco mais de um quarto das vendas de eletrificados,
mas por cerca de três quartos da massa de baterias vendidas.

\paragraph{Cenários de eletrificação (Objetivo~2).}
Os cenários Conservador e Acelerado, construídos a partir do estudo ANFAVEA/BCG,
delimitam uma frota eletrificada de veículos leves de 5--6~milhões de veículos em
2030, 12--14~milhões em 2035 e 26--28~milhões em 2040, incluídos os híbridos leves;
os BEV e PHEV, que concentram a massa de baterias, passam de 2--3~milhões para
11--14~milhões. A faixa é intermediária entre as projeções disponíveis: a do
PDE~2035 é muito inferior, a da Bright Consulting para 2030 é superior, e a do MBC
Brasil para 2040 é próxima do cenário Conservador sem os híbridos leves.

\paragraph{Fluxos da fabricação de células LFP (Objetivo~3).}
No caso central, a fabricação de 1~kWh de célula LFP (6,25~kg) consome 6,66~kg de
materiais e gera 0,39~kg de refugo, com rendimento global de massa de 93,9\%; 97\%
do refugo é destinado à valorização externa, e a pasta descartada (3\%), à
disposição como resíduo perigoso. O NMP opera em circuito praticamente fechado
(1,69 dos 1,70~kg/kWh evaporados são recuperados), e o consumo de eletricidade
(35~kWh/kWh) concentra-se na secagem e na sala seca (64\%), operações que não
transformam o material — o que explica o interesse no revestimento a seco.
Normalizado para 1~GWh/ano, o sistema consome cerca de 6.660~t de materiais e
35~GWh de eletricidade e gera cerca de 390~t de refugo. A ampliação do fluxograma
mostrou que as operações anteriores e posteriores às etapas de fabricação não
alteram o balanço de massa em regime permanente, mas concentram funções críticas
— controle de umidade dos insumos, segregação segura das células reprovadas e
rastreabilidade —, ainda não quantificadas. Os resultados são sensíveis às taxas
de refugo por etapa: com valores intermediários das faixas da literatura, o refugo
quase dobra.

\paragraph{Implicações para a cadeia de reciclagem (Objetivo~4).}
A sucata de produção — homogênea, de composição conhecida e disponível desde o
início da operação — é o primeiro fluxo de material secundário de uma eventual
cadeia nacional de reciclagem e é compatível com a reciclagem direta, rota mais
promissora para o LFP. Já a reciclagem das baterias pós-consumo de LFP permanece
condicionada à escala, ao preço do lítio e à construção de um marco regulatório
específico, hoje em discussão no Senado Federal (PL~2.132/2025).

\paragraph{Limitações.}
O \textit{market share} foi calculado em número de veículos, sobre uma base cujo
período e denominadores ainda precisam ser documentados; os cenários dependem de
um único estudo como espinha dorsal, com valores arredondados ao milhão; e a MEFA
baseia-se em faixas da literatura, normalizadas para uma planta madura, com
operações de suporte apenas qualitativas e fatores de emissão aproximados. As
pendências de fonte e de dado estão sinalizadas ao longo do texto.

\paragraph{Continuidade do trabalho.}
A continuidade do trabalho compreende: (i)~o MFA dinâmico da frota — estoques e
fluxos de baterias em uso, em fim de vida e em segunda vida, para lítio, cobalto,
níquel, manganês, ferro, fósforo, grafite e, no caso das baterias NiMH, elementos
de terras raras —, segundo a metodologia de \citeonline{BrunnerRechberger} e
implementado no \textit{software} STAN; (ii)~a conversão dos cenários de frota em
demanda anual de células, com capacidades médias de bateria por categoria;
(iii)~a quantificação das operações de suporte da fábrica e a caracterização das
rotas de calor e de água; e (iv)~a avaliação dos impactos ambientais e sociais da
cadeia de baterias, tema do projeto ainda não abordado nesta etapa.
